% Part: normal-modal-logic
% Chapter: filtrations
% Section: S5-decidable

\documentclass[../../../include/open-logic-section]{subfiles}

\begin{document}

\olfileid{nml}{fil}{dec}

\olsection{\Log{S5} નિર્ણેય છે}
સાંત નિદર્શન ગુણધર્મથી સૂત્ર-ઢાંચાઓ દ્વારા અપાયેલી નિદર્શ્ય-તર્ક
પદ્ધતિઓ \emph{નિર્ણેય} છે તે બતાવવાનો સરળ માર્ગ મળે છે (એટલે કે
કોઈ !!a{formula} પદ્ધતિમાં !!{derivable} છે કે નહિ તે નક્કી કરવાની
ગણનક્ષમ પદ્ધતિ છે).

\begin{thm}
  \Log{S5} નિર્ણેય છે.
\end{thm}

\begin{proof}
  $!A$ આપેલું હોય અને $!A$માં આવતાં !!{propositional variable}s
  $p_1$, \dots, $p_k$માંનાં હોય એમ ધારો. દરેક $n$ માટે $n$ જગતવાળાં
  અને $p_1$, \dots, $p_k$ને સત્યમૂલ્ય આપતાં સાંત સંખ્યામાં જ નિદર્શનો
  હોવાથી આપણે બે પ્રક્રિયાઓ \emph{સમાંતરે} ચલાવી શકીએ: કોઈ વ્યવસ્થિત
  રીતે પ્રમાણો રચીને \Log{S5}નાં બધાં પ્રમેયોની ગણના કરવી; અને
  $1$, $2$, \dots જગતવાળાં બધાં સાર્વત્રિક નિદર્શનોની ગણના કરીને
  તેમાંનાં કોઈ જગત પર $!A$ અસત્ય છે કે નહિ તે ચકાસવું.
  \sourcecorrection{OLMD-047}{મૂળમાં બીજી પ્રક્રિયા માટે બધાં સાંત
  નિદર્શનોની ગણના લખી છે; તે S5નું નિર્ણેયન આપતું નથી, કારણ કે
  અસર્વત્રિક નિદર્શન S5-પ્રમેયને પણ ખોટું પાડી શકે. S5ના
  સાર્વત્રિક નિદર્શનો સુધી ગણના મર્યાદિત કરી છે.}
  અંતે બેમાંથી એક પ્રક્રિયા ઉત્તર આપશે: કારણ કે
  \olref[com][fra]{thm:generaldet} અને \olref[fmp]{cor:S5fmp} મુજબ,
  કાં તો $!A$ !!{derivable} છે, અથવા તે કોઈ સાંત સાર્વત્રિક
  નિદર્શનમાં અસત્ય છે.
\end{proof}

ઉપરની સાબિતી \Log{S5} માટે કામ કરે છે કારણ કે સાર્વત્રિક
નિદર્શનોનાં ફિલ્ટ્રેશનો આપમેળે સાર્વત્રિક હોય છે. સ્વવાચકતા અને
સિરિયલતા માટે પણ એવું જ છે, પરંતુ અન્ય ગુણધર્મો માટે વધુ કામ
જરૂરી છે.

\end{document}
